\documentclass[12pt]{article}
\usepackage[utf8]{inputenc}
\usepackage{amsmath}
\usepackage{amssymb}
\usepackage{geometry}
\usepackage{listings}
\usepackage{xcolor}
\usepackage{hyperref}
\usepackage{booktabs}
\usepackage{array}
\usepackage{graphicx}
\usepackage{float}
\usepackage{enumitem}

\geometry{a4paper, margin=1in}

\definecolor{codegreen}{rgb}{0,0.6,0}
\definecolor{codegray}{rgb}{0.5,0.5,0.5}
\definecolor{codepurple}{rgb}{0.58,0,0.82}
\definecolor{backcolour}{rgb}{0.95,0.95,0.92}

\lstset{
    backgroundcolor=\color{backcolour},
    commentstyle=\color{codegreen},
    keywordstyle=\color{magenta},
    numberstyle=\tiny\color{codegray},
    stringstyle=\color{codepurple},
    basicstyle=\ttfamily\small,
    breakatwhitespace=false,
    breaklines=true,
    captionpos=b,
    keepspaces=true,
    numbers=left,
    numbersep=5pt,
    showspaces=false,
    showstringspaces=false,
    showtabs=false,
    tabsize=2,
    language=Python
}

\title{\textbf{Documentation: 3D Drone Rescue Navigation} \\
       \large A* Search Algorithm for Smart City Rescue Mission \\
       \large AI Lab — Assignment 4}
\author{Uttam Mahata (2022CSB104)}
\date{\today}

\begin{document}

\maketitle
\tableofcontents
\newpage

% ─────────────────────────────────────────────────────────────────────────────
\section{Introduction}
% ─────────────────────────────────────────────────────────────────────────────

A disaster has struck a futuristic smart city.
An AI-powered drone must autonomously navigate through the city, avoiding
collapsed buildings and fire zones, rescue all stranded survivors, and return
to base — while minimising total energy consumption.

This implementation uses the \textbf{A* search algorithm} extended to a
three-dimensional grid.  The drone operates in a layered city space (Ground,
Mid, and Sky levels) and must visit multiple intermediate goal states
(survivor locations) before reaching the final goal (base station).

Key challenges addressed:
\begin{itemize}
    \item 3D state-space navigation with heterogeneous movement costs.
    \item Multi-goal search using greedy nearest-neighbour ordering.
    \item Dynamic recharge logic when energy becomes insufficient mid-mission.
    \item Admissible 3D Euclidean heuristic that guarantees optimal paths per segment.
\end{itemize}

% ─────────────────────────────────────────────────────────────────────────────
\section{Problem Formulation}
% ─────────────────────────────────────────────────────────────────────────────

\subsection{State Space}

The smart city is modelled as a 3-dimensional grid.
Each cell is identified by coordinates $(x, y, z)$ where:
\begin{itemize}
    \item $x \in \{0, \ldots, \text{rows}-1\}$ — row index (North–South axis).
    \item $y \in \{0, \ldots, \text{cols}-1\}$ — column index (East–West axis).
    \item $z \in \{0, \ldots, \text{levels}-1\}$ — altitude level (0 = Ground, 1 = Mid, 2 = Sky).
\end{itemize}

\subsection{Cell Types}

\begin{table}[H]
\centering
\begin{tabular}{@{} c l l @{}}
\toprule
\textbf{Symbol} & \textbf{Type} & \textbf{Description} \\
\midrule
\texttt{0}  & Road / Open space     & Passable; base terrain cost = 1 \\
\texttt{1}  & Building / Obstacle   & Impassable; drone cannot pass \\
\texttt{F}  & Fire zone             & Passable; high-risk; terrain cost = 3 \\
\texttt{S}  & Drone start           & Passable (cost = 1); mission origin \\
\texttt{Gn} & Survivor location     & Passable (cost = 1); must be visited \\
\texttt{R}  & Recharging station    & Passable (cost = 1); restores energy \\
\texttt{B}  & Base station          & Passable (cost = 1); final goal \\
\bottomrule
\end{tabular}
\caption{Cell types and their properties.}
\end{table}

\subsection{Initial and Goal States}

\begin{itemize}
    \item \textbf{Initial state:} the drone is at position $S$ with energy $E_0$.
    \item \textbf{Intermediate goals:} all survivor locations $G_1, G_2, \ldots, G_n$
          must be visited (in any order).
    \item \textbf{Final goal:} the drone must return to base $B$ after rescuing
          all survivors.
\end{itemize}

% ─────────────────────────────────────────────────────────────────────────────
\section{Movement Model}
% ─────────────────────────────────────────────────────────────────────────────

\subsection{Six-Directional Movement}

The drone moves in six axis-aligned directions — no diagonal movement is
permitted:

\begin{center}
\begin{tabular}{@{} l l @{}}
\toprule
\textbf{Direction} & \textbf{Coordinate change} \\
\midrule
North  & $(x-1,\; y,\; z)$ \\
South  & $(x+1,\; y,\; z)$ \\
East   & $(x,\; y+1,\; z)$ \\
West   & $(x,\; y-1,\; z)$ \\
Up     & $(x,\; y,\; z+1)$ \\
Down   & $(x,\; y,\; z-1)$ \\
\bottomrule
\end{tabular}
\end{center}

\subsection{Energy Cost Model}

The energy cost $c(u, v)$ to move from cell $u$ to adjacent cell $v$ depends
on both the \emph{terrain} of the destination cell and the \emph{direction} of
movement:

\begin{equation}
c(u, v) =
\begin{cases}
\tau(v)       & \text{if horizontal move } (\Delta z = 0) \\
\tau(v) + 2   & \text{if ascending } (\Delta z = +1) \\
1             & \text{if descending } (\Delta z = -1)
\end{cases}
\end{equation}

where the terrain cost $\tau(v)$ is:

\begin{equation}
\tau(v) =
\begin{cases}
3 & \text{if } v \text{ is a fire zone (F)} \\
1 & \text{otherwise (open space, S, B, R, or } G_n\text{)}
\end{cases}
\end{equation}

\textbf{Rationale:}
\begin{itemize}
    \item Horizontal flight through free space costs 1 unit.
    \item Fire zones create turbulence — horizontal traversal costs 3 units.
    \item Ascending requires extra thrust against gravity: $+2$ energy units on
          top of the terrain cost.
    \item Descending exploits gravity assist and always costs just 1 unit,
          regardless of terrain below.
\end{itemize}

% ─────────────────────────────────────────────────────────────────────────────
\section{A* Algorithm}
% ─────────────────────────────────────────────────────────────────────────────

\subsection{Evaluation Function}

A* evaluates each node $n$ using the combined function:
\begin{equation}
    f(n) = g(n) + h(n)
\end{equation}

\begin{itemize}
    \item $g(n)$ — actual energy consumed from the start to node $n$.
    \item $h(n)$ — admissible heuristic estimate of the energy from $n$ to the goal.
    \item $f(n)$ — estimated total energy for the best path through $n$.
\end{itemize}

\subsection{Heuristic Function}

The heuristic uses the \textbf{3D Euclidean distance}:

\begin{equation}
h(n) = \sqrt{(x_n - x_g)^2 + (y_n - y_g)^2 + (z_n - z_g)^2}
\end{equation}

\textbf{Admissibility:}
Since the minimum possible energy for any single step is 1 unit (descending
through free space), and the Euclidean distance never overestimates the
number of steps required ($h_{\text{Euclidean}} \leq d_{\text{L1}}$), we have:
\begin{equation}
h(n) \leq h^*(n) \quad \forall n
\end{equation}
where $h^*(n)$ is the true minimum energy cost to the goal.
This guarantees A* finds the \emph{optimal} (minimum energy) path per segment.

\subsection{Priority Queue}

Nodes are stored in a \textbf{min-heap} (priority queue) ordered by $f(n)$.
A tie-breaking counter ensures that no two heap entries are compared on their
position tuples:

\begin{lstlisting}[caption={Priority queue push with tie-breaking.}]
heapq.heappush(open_set, (f_score, g_score, next(counter), position))
\end{lstlisting}

Stale (outdated) entries are lazily discarded:

\begin{lstlisting}[caption={Skipping outdated entries.}]
if g_val > g_score.get(current, float('inf')):
    continue   # outdated entry — skip
\end{lstlisting}

\subsection{Algorithm Pseudocode}

\begin{lstlisting}[language={}, caption={A* on 3D City Grid (pseudocode)}, mathescape=true]
Input:  city grid G, start s, goal t
Output: shortest path P and cost c,  or None if unreachable

  open  <- { (h(s,t), 0, s) }          -- min-heap: (f, g, node)
  g[s]  <- 0 ;  g[v] <- inf  for all v != s
  came_from <- {}

  while open != empty:
      (f, g_val, u) <- open.pop_min()
      if g_val > g[u]:  continue        -- stale entry
      if u = t:  return reconstruct(came_from, u), g[t]
      for each neighbour v of u:
          g' <- g[u] + c(u, v)
          if g' < g[v]:
              g[v]          <- g'
              came_from[v]  <- u
              open.push( g' + h(v,t), g', v )

  return None, inf
\end{lstlisting}

% ─────────────────────────────────────────────────────────────────────────────
\section{Multi-Goal Search Strategy}
% ─────────────────────────────────────────────────────────────────────────────

\subsection{Segment-by-Segment A*}

The full mission is decomposed into independent A* searches:
\[
S \xrightarrow{A^*} G_{\pi(1)} \xrightarrow{A^*} G_{\pi(2)} \xrightarrow{A^*}
\cdots \xrightarrow{A^*} G_{\pi(n)} \xrightarrow{A^*} B
\]
where $\pi$ is the visit permutation determined by the goal ordering heuristic.

\subsection{Greedy Nearest-Neighbour Goal Ordering}

Finding the optimal visit order for $n$ survivors is equivalent to the
Travelling Salesman Problem (NP-hard).  For practical sizes, a
\textbf{greedy nearest-neighbour} heuristic is applied:

\begin{lstlisting}[language={}, caption={Greedy Nearest-Neighbour Goal Ordering (pseudocode)}]
Input:  start s,  goal set G
Output: ordered list of goals

  ordered <- [] ;  current <- s ;  R <- G
  while R != empty:
      g* <- argmin_{g in R}  h(current, g)
      ordered.append(g*)
      current <- g*
      R <- R \ {g*}
  return ordered
\end{lstlisting}

The Euclidean distance $h$ is used for proximity estimation.

% ─────────────────────────────────────────────────────────────────────────────
\section{Recharge Logic}
% ─────────────────────────────────────────────────────────────────────────────

\subsection{Energy Tracking and Recharge Trigger}

The drone tracks remaining energy $E$ throughout the mission.
Before each segment, the A* cost $c_{\text{seg}}$ to the next goal is computed.
If $c_{\text{seg}} > E$, the recharge procedure is triggered:

\begin{enumerate}
    \item Find the nearest reachable recharging station $R^*$ via A*.
    \item If $\text{cost}(R^*) > E$: \textbf{Mission failure} — energy depleted
          before reaching any station.
    \item Fly to $R^*$, deducting travel cost from $E$.
    \item \textbf{Recharge:} $E \gets E_0 - p$, where $E_0$ is the initial energy
          and $p$ is the configurable recharge penalty.
    \item Re-run A* from $R^*$ to the original target goal.
\end{enumerate}

\subsection{Recharge Penalty}

Stopping at a recharge station incurs a fixed energy penalty $p$ (default: 5
units) to model the time cost of landing, docking, and recharging.
After recharging, the drone has energy $E_0 - p$ available (full recharge
minus the penalty), ensuring the penalty discourages unnecessary detours.

% ─────────────────────────────────────────────────────────────────────────────
\section{Implementation}
% ─────────────────────────────────────────────────────────────────────────────

\subsection{Module Structure}

\begin{lstlisting}[caption={Module-level layout of drone\_navigation.py}]
drone_navigation.py
├── Constants  (OBSTACLE, FIRE, DIRECTIONS, LEVEL_NAMES)
├── CityGrid                    # loads and queries the 3D grid
│   ├── __init__(data)          # parse JSON, locate S/B/Gn/R
│   ├── get_cell(x, y, z)       # bounds-safe cell lookup
│   ├── terrain_cost(cell)      # 3 for F, 1 otherwise
│   ├── is_passable(cell)       # False for obstacle '1'
│   ├── move_cost(dest, dz)     # direction-aware energy cost
│   └── get_neighbors(pos)      # valid (neighbour, cost) pairs
├── heuristic(a, b)             # 3D Euclidean distance
├── astar_3d(grid, start, goal) # core A* — returns (path, cost)
├── _reconstruct_path(...)      # backtrack came_from
├── _greedy_goal_order(...)     # nearest-neighbour TSP heuristic
├── _find_nearest_recharge(...) # A* to closest R station
├── plan_mission(...)           # orchestrates full mission
├── visualize_grid(...)         # ASCII layer-by-layer output
└── main()                      # CLI entry point
\end{lstlisting}

\subsection{CityGrid Class}

The \texttt{CityGrid} class wraps the JSON grid data and provides
neighbour-lookup with full cost computation:

\begin{lstlisting}[caption={CityGrid.get\_neighbors — cost-aware 3D neighbourhood.}]
def get_neighbors(self, pos):
    x, y, z = pos
    result = []
    for dx, dy, dz, _ in DIRECTIONS:
        nx, ny, nz = x + dx, y + dy, z + dz
        cell = self.get_cell(nx, ny, nz)
        if self.is_passable(cell):
            result.append(((nx, ny, nz), self.move_cost(cell, dz)))
    return result
\end{lstlisting}

\subsection{A* Core}

The \texttt{astar\_3d} function uses a min-heap with lazy deletion of stale entries:

\begin{lstlisting}[caption={Core A* loop.}]
while open_set:
    f_val, g_val, _, current = heapq.heappop(open_set)

    if current == goal:
        return _reconstruct_path(came_from, current), g_score[goal]

    if g_val > g_score.get(current, float('inf')):
        continue   # stale entry

    for neighbor, cost in city_grid.get_neighbors(current):
        tentative = g_score[current] + cost
        if tentative < g_score.get(neighbor, float('inf')):
            came_from[neighbor] = current
            g_score[neighbor]   = tentative
            f_new = tentative + heuristic(neighbor, goal)
            heapq.heappush(open_set,
                           (f_new, tentative, next(_TIE_COUNTER), neighbor))
\end{lstlisting}

\subsection{Input Format}

The city grid is provided as a \textbf{JSON file} with the following schema:

\begin{lstlisting}[language={},caption={city\_grid.json structure.}]
{
  "dimensions": { "rows": 5, "cols": 5, "levels": 3 },
  "initial_energy": 100,
  "recharge_penalty": 5,
  "grid": [
    [ /* z=0 Ground Level */ [...], [...], ... ],
    [ /* z=1 Mid Level    */ [...], [...], ... ],
    [ /* z=2 Sky Level    */ [...], [...], ... ]
  ]
}
\end{lstlisting}

\subsection{Command-Line Interface}

\begin{lstlisting}[language=bash,caption={Usage.}]
python drone_navigation.py city_grid.json
python drone_navigation.py city_grid.json --energy 80 --recharge-penalty 10
\end{lstlisting}

\begin{center}
\begin{tabular}{@{} l l l @{}}
\toprule
\textbf{Argument}          & \textbf{Default}     & \textbf{Description} \\
\midrule
\texttt{grid\_file}         & —                    & Path to city grid JSON \\
\texttt{--energy}           & from JSON (100)      & Initial drone energy \\
\texttt{--recharge-penalty} & from JSON (5)        & Energy penalty for recharging \\
\bottomrule
\end{tabular}
\end{center}

% ─────────────────────────────────────────────────────────────────────────────
\section{Example: 5×5×3 City Grid}
% ─────────────────────────────────────────────────────────────────────────────

\subsection{Grid Layout}

The example grid has 3 levels (z = 0, 1, 2) each of size 5×5.

\textbf{Ground Level (z = 0):}
\[
\begin{array}{|c|c|c|c|c|}
\hline
S  & 0  & \mathbf{1} & 0  & 0 \\
\hline
0  & F  & \mathbf{1} & \mathbf{1} & G_1 \\
\hline
F  & 0  & 0  & 0  & F \\
\hline
R  & \mathbf{1} & \mathbf{1} & \mathbf{1} & 0 \\
\hline
B  & 0  & \mathbf{1} & 0  & G_2 \\
\hline
\end{array}
\]

\textbf{Mid Level (z = 1):}
\[
\begin{array}{|c|c|c|c|c|}
\hline
0  & 0  & \mathbf{1} & 0  & 0 \\
\hline
0  & \mathbf{1} & \mathbf{1} & 0  & 0 \\
\hline
F  & 0  & 0  & 0  & \mathbf{1} \\
\hline
0  & \mathbf{1} & \mathbf{1} & \mathbf{1} & 0 \\
\hline
0  & 0  & 0  & 0  & 0 \\
\hline
\end{array}
\]

\textbf{Sky Level (z = 2):}
\[
\begin{array}{|c|c|c|c|c|}
\hline
\mathbf{1} & 0  & 0  & 0  & 0 \\
\hline
0  & F  & 0  & \mathbf{1} & 0 \\
\hline
0  & 0  & \mathbf{1} & 0  & F \\
\hline
0  & 0  & 0  & 0  & 0 \\
\hline
0  & \mathbf{1} & 0  & 0  & \mathbf{1} \\
\hline
\end{array}
\]

\subsection{Special Positions}

\begin{center}
\begin{tabular}{@{} l l @{}}
\toprule
\textbf{Entity} & \textbf{Position $(x, y, z)$} \\
\midrule
Start $S$         & $(0, 0, 0)$ \\
Survivor $G_1$    & $(1, 4, 0)$ \\
Survivor $G_2$    & $(4, 4, 0)$ \\
Recharge $R$      & $(3, 0, 0)$ \\
Base $B$          & $(4, 0, 0)$ \\
\bottomrule
\end{tabular}
\end{center}

% ─────────────────────────────────────────────────────────────────────────────
\section{Edge Cases}
% ─────────────────────────────────────────────────────────────────────────────

\begin{enumerate}
    \item \textbf{No valid path to a survivor:}
          A* returns \texttt{None} for that segment.
          The mission halts with an error message.

    \item \textbf{Energy depleted before reaching a recharge station:}
          If $\text{cost}(R^*) > E$, the drone cannot reach any station.
          The mission is declared failed.

    \item \textbf{Base station unreachable:}
          A* on the final segment returns \texttt{None}.
          Error reported; partial path and consumed energy are still output.

    \item \textbf{No survivor goals in grid:}
          The drone flies directly from $S$ to $B$ in a single A* segment.

    \item \textbf{No recharge stations in grid:}
          If energy runs out mid-mission, mission fails with a clear error
          rather than an infinite loop.

    \item \textbf{Start equals goal:}
          \texttt{astar\_3d} short-circuits and returns \texttt{([start], 0.0)}.
\end{enumerate}

% ─────────────────────────────────────────────────────────────────────────────
\section{Complexity Analysis}
% ─────────────────────────────────────────────────────────────────────────────

Let $N = \text{rows} \times \text{cols} \times \text{levels}$ be the total
number of cells and $k$ be the number of survivors.

\begin{table}[H]
\centering
\begin{tabular}{@{} l l l @{}}
\toprule
\textbf{Component}           & \textbf{Time}         & \textbf{Space} \\
\midrule
Single A* call               & $O(N \log N)$         & $O(N)$ \\
Goal ordering (greedy)       & $O(k^2)$              & $O(k)$ \\
Find nearest recharge        & $O(|\mathcal{R}| \cdot N \log N)$ & $O(N)$ \\
Full mission ($k+1$ segments)& $O(k \cdot N \log N)$ & $O(N)$ \\
\bottomrule
\end{tabular}
\caption{Asymptotic complexity ($\mathcal{R}$ = set of recharge stations).}
\end{table}

The space complexity per A* call is $O(N)$ for the \texttt{g\_score} dict and
the \texttt{came\_from} backtracking map.

% ─────────────────────────────────────────────────────────────────────────────
\section{Conclusion}
% ─────────────────────────────────────────────────────────────────────────────

This implementation extends the classical A* algorithm to a fully 3D,
heterogeneous-cost, multi-goal rescue scenario.
Key design decisions:

\begin{itemize}
    \item \textbf{Admissible 3D heuristic} — 3D Euclidean distance guarantees
          optimality for each individual segment.
    \item \textbf{Direction-aware costs} — ascending and descending have
          physically motivated asymmetric energy penalties.
    \item \textbf{Greedy goal ordering} — practical $O(k^2)$ nearest-neighbour
          heuristic avoids NP-hard TSP exact search.
    \item \textbf{Lazy recharge} — recharging is only triggered when necessary,
          minimising the fixed penalty cost.
    \item \textbf{Configurable parameters} — initial energy and recharge penalty
          can be specified via JSON or CLI flags, enabling scenario testing.
\end{itemize}

\end{document}
